\documentclass[11pt]{article}
\usepackage[margin=1in]{geometry}
\usepackage{amsmath,amssymb,amsthm,mathtools}
\usepackage[T1]{fontenc}
\usepackage{lmodern}
\usepackage{microtype}
\usepackage[colorlinks=true,linkcolor=blue,urlcolor=blue,citecolor=blue]{hyperref}
\newtheorem{theorem}{Theorem}
\newtheorem{lemma}{Lemma}
\newtheorem{remark}{Remark}
\newcommand{\Herm}{\operatorname{Herm}}
\newcommand{\Tr}{\operatorname{Tr}}
\newcommand{\Sch}{\operatorname{Sch}}
\newcommand{\supp}{\operatorname{supp}}
\newcommand{\Z}{\mathbb Z}
\newcommand{\R}{\mathbb R}
\newcommand{\C}{\mathbb C}
\newcommand{\Bcal}{\mathcal B}
\newcommand{\one}{\mathbf 1}
\newcommand{\norm}[1]{\left\|#1\right\|}
\newcommand{\abs}[1]{\left|#1\right|}
\title{Toward a four-dimensional Yang--Mills construction\\
\large Notebook 04: strictly local Gaussian cutoff activities}
\author{Research working document}
\date{16 October 2026}
\begin{document}
\maketitle

\noindent\textbf{Status: internally checked abstract theorem.}
This standalone notebook records a sufficient theorem and proof. The support
algebra, inverse estimate, full Gaussian density majorant, forest expansion,
and final support-root summation have received two substantive internal
reviews, including an independent derivation/audit. This is not external
mathematical certification. The hypotheses have not been matched to the
actual Yang--Mills fluctuation operators, and no interacting cutoff-transfer
theorem or continuum construction is claimed.

\medskip
\noindent
The purpose is to upgrade spatial localization in the positions of rare
Gaussian cutoff defects to strict locality in the background variables.
The extra hypothesis is a small sum of local perturbations of a fixed
precision matrix. The proof expands the entire normalized marginal density,
including its determinant and inverse. No derivative of a sharp indicator
is taken.

\section{Setting, constants, and conclusion}

Let $\Lambda\subset\Z^d$ be finite. Distances are nearest-neighbor lattice
graph distances, and each block carries $D$ real coordinates, with $D$ fixed.
Every constant below is uniform in $\Lambda$. The background domain is a
product
\[
 \mathcal U_\Lambda=\prod_{b\in\Lambda}\mathcal U_b
\]
of open complex domains. A local function with support $Y$ is a bounded
holomorphic function of $A|_Y$ on $\prod_{b\in Y}\mathcal U_b$.
Its uniform norm is denoted $\norm{\cdot}_\infty$.

Let $C_{0,\Lambda}=Q_{0,\Lambda}^{-1}$ be background-independent, complex
symmetric, and satisfy
\begin{equation}
 \Herm C_{0,\Lambda}\ge\lambda I,\qquad
 \norm{C_{0,\Lambda}}\le\Lambda_0,
 \qquad 0<\lambda\le\Lambda_0<\infty.
 \label{eq:baseline-gap}
\end{equation}
Real positive $C_{0,\Lambda}$ is the principal example. Fix $0<\kappa<m$
such that the weighted block Schur bound
\begin{equation}
 K:=\sup_\Lambda\sup_{b\in\Lambda}
       \sum_{c\in\Lambda}
       e^{(\kappa+2m)d(b,c)}\norm{(C_{0,\Lambda})_{bc}}
       <\infty
 \label{eq:baseline-decay}
\end{equation}
holds. Transpose symmetry gives the same column bound.

Consider the precision
\begin{equation}
 Q_\Lambda(A)=Q_{0,\Lambda}+V_\Lambda(A),\qquad
 V_\Lambda(A)=\sum_{\varnothing\ne Y\subset\Lambda}V_Y(A|_Y),
 \label{eq:precision}
\end{equation}
where each $V_Y$ is individually complex symmetric, is supported in
$Y\times Y$ as a block matrix, and is holomorphic and local as above.
Write $\ell(Y)$ for the minimum metric spanning-tree length of $Y$,
with $\ell(\{b\})=0$. Assume
\begin{equation}
 \nu:=\sup_\Lambda\sup_{b\in\Lambda}
       \sum_{Y\ni b}\sum_{c\in Y}
       e^{(\kappa+2m)\ell(Y)}\norm{(V_Y)_{bc}}_\infty
       <\infty.
 \label{eq:local-norm}
\end{equation}
The same bound holds for columns. In particular the ordinary Schur norm
of $V_\Lambda$ is at most $\nu$.

Here are explicit sufficient smallness conditions. Define
\begin{align}
 \delta&=\lambda/4,&
 v&=\frac{K^2\nu}{1-K\nu},&
 v_1&=v(1+\delta/K),                                      \label{eq:vconstants}\\
 s&=\min\left\{\frac m2,\frac{\lambda e m}{8(K+\delta)}\right\},&
 L&=\frac4\lambda\sum_{z\in\Z^d}e^{-s|z|},&
 V_*&=\frac{8\Lambda_0^2}{\lambda}.                       \label{eq:inverseconstants}
\end{align}
Require
\begin{equation}
 K\nu<1,\qquad v_1\le\lambda/4,\qquad
 u:=Lv_1\le\frac12,\qquad
 w:=\frac{L^2v_1}{1-u}\le\frac1{2V_*}.
 \label{eq:smallness}
\end{equation}
These conditions hold when $\nu$ is sufficiently small with the baseline
constants fixed. Set
\begin{align}
 \widehat\tau_R
   &=\left(\frac{16V_*}{\lambda}\right)^{D/2}
                         e^{-R^2/(8V_*)},                 \label{eq:tau}\\
 A_\kappa
   &=K\sum_{z\in\Z^d\setminus\{0\}}e^{-(m-\kappa)|z|},
 &B_0&=\frac{8e^2A_\kappa}{\delta},
 &a_*&=\min\{1/2,\kappa/4\}.                              \label{eq:clusterconstants}
\end{align}
Finally require
\begin{equation}
 \widehat\tau_R\le\frac1{2e^2\max\{1,B_0\}}.
 \label{eq:cutoff-smallness}
\end{equation}

For a complex symmetric accretive covariance $C$, use the normalized
Gaussian density
\begin{equation}
 p_C(x)=(2\pi)^{-N/2}(\det C)^{-1/2}
                           e^{-x^TC^{-1}x/2}.
 \label{eq:density}
\end{equation}
The determinant branch is the analytic branch on the accretive cone that
agrees with the positive square root on real positive matrices. This
density integrates to one. Its real-domain integral is absolutely
convergent because
\[
 \Herm C^{-1}=C^{-*}(\Herm C)C^{-1}>0.
\]
Put $C_\Lambda(A)=Q_\Lambda(A)^{-1}$ and define the fixed-domain cutoff
partition function
\begin{equation}
 Z_{\Lambda,R}(A)=
  \int_{\R^{D|\Lambda|}}p_{C_\Lambda(A)}(x)
        \prod_{b\in\Lambda}\one_{\{\norm{x_b}_2\le R\}}\,dx.
 \label{eq:partition}
\end{equation}
The real domain and threshold do not move with the complex background.

\begin{theorem}[Proposed strict-locality theorem]\label{thm:main}
Under \eqref{eq:baseline-gap}--\eqref{eq:cutoff-smallness},
$Z_{\Lambda,R}$ has a holomorphic logarithm, fixed by turning on the
defect strength from zero, and an absolutely convergent representation
\begin{equation}
 \log Z_{\Lambda,R}(A)
       =\sum_{H}E_{\Lambda,H}(A|_{H\cap\Lambda}).
 \label{eq:local-log}
\end{equation}
Here $H$ ranges over finite connected sets in $\Z^d$, and the activities are
strictly local, bounded, and holomorphic in the indicated variable
background. They satisfy
\begin{equation}
 \sup_{b\in\Z^d}\sum_{H\ni b}
     e^{a_*|H|}\norm{E_{\Lambda,H}}_\infty
 \le
 2e^2\left(\sum_{z\in\Z^d}e^{-(\kappa/2)|z|}\right)
                     \widehat\tau_R.
 \label{eq:final-bound}
\end{equation}
One may take $H$ in the lattice hull of $\Lambda$ when $\Lambda$ is not
geodesically convex.
\end{theorem}

\begin{remark}[Finite-volume dependence]
The activity $E_{\Lambda,H}$ may depend on the entire
\emph{background-independent} baseline $C_{0,\Lambda}$. Strict locality
asserts only that its \emph{variable background dependence} reads
$A|_{H\cap\Lambda}$. Neither equality of activities in different volumes
nor existence of their thermodynamic limit follows from
\eqref{eq:final-bound}. A consistency or comparison argument for the
baselines is an additional obligation.
\end{remark}

\begin{remark}[Product-domain hypothesis]
The product background domain is an explicit sufficient hypothesis.
A weak gauge-potential ball or a constrained gauge domain is not
automatically such a product. An application must supply an appropriate
product chart or polydomain, or separately verify the support algebra
and local-function bounds on the actual domain.
\end{remark}

\section{The support algebra and local covariance paths}

For the proof, suppress $\Lambda$ from the notation. Fix a nonempty
marked set $S\subset\Lambda$. For a background support $X\subset\Lambda$,
let $\ell_S(X)$ be the minimum total length of a metric forest on
$S\cup X$ in which every component meets $S$.
Roots in $S$ may remain in different components, and
$\ell_S(\varnothing)=0$. We have
\begin{equation}
 \ell_S(X\cup X')\le\ell_S(X)+\ell_S(X'),\qquad
 S\subset S'\Longrightarrow\ell_{S'}(X)\le\ell_S(X).
 \label{eq:forest-subadd}
\end{equation}

Consider formal families $F=(F_X)_X$ of bounded local holomorphic
functions with norm
\begin{equation}
 \norm{F}_{\Bcal_S}
       =\sum_Xe^{\kappa\ell_S(X)}\norm{F_X}_\infty.
 \label{eq:supportnorm}
\end{equation}
Multiplication is union convolution,
$(FG)_Z=\sum_{X\cup X'=Z}F_XG_{X'}$.
Equation \eqref{eq:forest-subadd} makes the completion a commutative
unital Banach algebra. Evaluation
$F\mapsto\sum_XF_X(A|_X)$ is a bounded algebra homomorphism.
Representations need not be unique; throughout use the explicit series
in the $V_Y$, grouped by the union of their supports. Formal scalar
variables multiplying the $V_Y$ can be retained until absolute
convergence has been established.

For matrices with entries in this support algebra, use the maximum of
the row and column Schur norms, taking the block operator norm inside
the support sum. This norm is submultiplicative.

The covariance is the convergent Neumann series
\begin{equation}
 C(A)=\sum_{j\ge0}(-1)^j C_0(VC_0)^j.
 \label{eq:neumann}
\end{equation}
Let $\Delta C=C-C_0$, and denote by $\Delta C_S$ its actual principal
marginal on $S$.

\begin{lemma}[Covariance paths with attached supports]\label{lem:path}
The series \eqref{eq:neumann}, with its explicit local supports, obeys
\begin{equation}
 \max\left\{
 \sup_{b\in S}\sum_{c\in S}e^{2m d(b,c)}
                           \norm{(\Delta C_S)_{bc}}_{\Bcal_S},
 \sup_{c\in S}\sum_{b\in S}e^{2m d(b,c)}
                           \norm{(\Delta C_S)_{bc}}_{\Bcal_S}
 \right\}\le v.
 \label{eq:cov-supportbound}
\end{equation}
\end{lemma}

\begin{proof}
In a term
$C_0V_{Y_1}C_0\cdots V_{Y_j}C_0$ with endpoints $b,c\in S$,
choose a spanning tree for each $Y_i$. Together with the covariance
links between consecutive perturbations these trees connect all
background supports to the endpoints. If $L_{\rm path}$ is the sum
of the covariance-link distances and $\sum_i\ell(Y_i)$, then
\[
 \ell_S\left(\bigcup_iY_i\right)\le L_{\rm path},
 \qquad d(b,c)\le L_{\rm path}.
\]
Therefore
\[
 e^{\kappa\ell_S(\cup_iY_i)+2m d(b,c)}
          \le e^{(\kappa+2m)L_{\rm path}}.
\]
Successive weighted matrix row sums using
\eqref{eq:baseline-decay} and \eqref{eq:local-norm} bound the order-$j$
contribution by $K^{j+1}\nu^j$.
The column calculation is identical. Summing $j\ge1$ gives
$v=K^2\nu/(1-K\nu)$.
\end{proof}

In particular,
$\norm{(\Delta C_S)_{bc}}_{\Bcal_S}\le v e^{-2m d(b,c)}$.
This pays for the support of the \emph{full covariance marginal}.
It is stronger than expanding only the covariance factors that will
appear on a marked tree.

\section{Tree interpolation and uniform inverse bounds}

For a partition $\pi$ of $S$, let $v_\pi$ be the zero-one mask that
retains covariance entries within each component and sets other entries
to zero. The matrix $C_{0,S}(v_\pi)$ is a direct sum of principal
submatrices, so it obeys \eqref{eq:baseline-gap}.
The same is true for every mask in the convex hull of the partition
masks.

For a forest $F$ and parameters $t_e\in[0,1]$, the BKAR mask is
\begin{equation}
 s^F_{bc}(t)=
 \begin{cases}
 \min_{e\in\operatorname{path}_F(b,c)}t_e,
                 &b,c\text{ in the same component},\\
 0,              &b,c\text{ in different components},
 \end{cases}
 \label{eq:bkar-mask}
\end{equation}
with diagonal entries equal to one. It belongs to that convex hull:
integrate, over $r\in[0,1]$, the partition mask of the subforest
retaining edges with $t_e\ge r$. This gives \eqref{eq:bkar-mask}.
The forest formula may be applied directly on this partition polytope;
no positivity assumption for every arbitrary coordinate-cube mask
is needed \cite{AR,BKARnotes}.

Fix a tree $T$ on $S$, $n=|S|\ge2$, with vertex degrees $d_b$.
Put
\begin{equation}
 J_{bc}=K e^{-m d(b,c)},\qquad
 r_{\{b,c\}}=
       \frac{\delta}{\max(d_b,d_c)J_{bc}}
                  \quad(\{b,c\}\in T).
 \label{eq:cauchyradii}
\end{equation}
At a BKAR point $s^T$, allow independent complex increments $z_e$
on its tree edges with $|z_e|<r_e$. Let $M_0$ be the resulting
masked baseline $C_{0,S}(s^T+z)$.
Since
\[
 \sum_{c:\{b,c\}\in T}\frac1{\max(d_b,d_c)}\le1,
\]
the perturbation of $C_{0,S}(s^T)$ has row and column norms at most
$\delta$. Consequently
\begin{equation}
 \Herm M_0\ge\lambda I/2,\qquad
 \norm{M_0}\le2\Lambda_0.
 \label{eq:deformedgap}
\end{equation}
Moreover, its Schur norm with entry weight $e^{m d(b,c)}$ is at most
$K+\delta$. For the perturbation part, this uses
\[
 \norm{(C_0)_{bc}}\le K e^{-(\kappa+2m)d(b,c)}
\]
in \eqref{eq:cauchyradii}.

\begin{lemma}[Elementary weighted-conjugation inverse bound]
\label{lem:inverse}
Suppose $\Herm M\ge aI$ and the weighted Schur norm of $M$ at rate
$m$ is at most $\overline K$. Let
\[
 s_0=\min\{m/2,aem/(4\overline K)\}.
\]
Then
\begin{equation}
 \norm{(M^{-1})_{bc}}\le(2/a)e^{-s_0d(b,c)}.
 \label{eq:inverse-decay}
\end{equation}
\end{lemma}

\begin{proof}
Fix $c$ and conjugate by the scalar block-diagonal weight
$D_{bb}=e^{s_0d(b,c)}$. The triangle inequality for the distance gives
\begin{align*}
 \norm{DMD^{-1}-M}_{\Sch}
 &\le \overline K\sup_{r\ge0}(e^{s_0r}-1)e^{-mr}\\
 &\le \frac{2s_0\overline K}{em}\le a/2.
\end{align*}
The middle inequality follows from
$e^{s_0r}-1\le s_0r e^{s_0r}$ and
$\sup_{r\ge0}r e^{-(m-s_0)r}\le2/(em)$.
As $\norm{M^{-1}}\le1/a$, a Neumann inverse yields
$\norm{(DMD^{-1})^{-1}}\le2/a$.
Its $(b,c)$ block is $e^{s_0d(b,c)}(M^{-1})_{bc}$,
which proves \eqref{eq:inverse-decay}.
\end{proof}

Apply Lemma \ref{lem:inverse} with $a=\lambda/2$ and
$\overline K=K+\delta$. The resulting rate is the $s$ in
\eqref{eq:inverseconstants}, and
\begin{equation}
 \norm{M_0^{-1}}_{\Sch}\le L.
 \label{eq:inverse-schur}
\end{equation}
This is uniform in $S,T,t,z$. The rate $s$ may be much smaller than
$m$ or $\kappa$. Only the \emph{ordinary} inverse Schur norm $L$ is
used below. All its matrix indices lie in the free root set $S$;
its entries do not pay for attachment of background supports.
Thus the proof does not require $\kappa<s$.

\section{Expansion of the entire marginal Gaussian density}

Apply the same mask $s^T+z$ to $C_S(A)=C_{0,S}+\Delta C_S$.
Write the resulting matrix as $M=M_0+E$. Lemma \ref{lem:path}
and \eqref{eq:cauchyradii} imply
\begin{equation}
 \norm{E}_{\Sch(\Bcal_S)}\le v_1=v(1+\delta/K).
 \label{eq:maskedE}
\end{equation}
The $s^T$ part has norm at most $v$.
The additional tree-edge row sums are bounded by
$\delta v/K$, using the entry estimate
$v e^{-2m d(b,c)}$ and the degree sum above.
The column bound is the same.

In the matrix support algebra set $F=M_0^{-1}E$.
By \eqref{eq:smallness}, $\norm{F}_{\Sch(\Bcal_S)}\le u\le1/2$,
so the full inverse and determinant expansions converge:
\begin{align}
 M^{-1}-M_0^{-1}
    &=\sum_{j\ge1}(-F)^jM_0^{-1},&
 \norm{M^{-1}-M_0^{-1}}_{\Sch(\Bcal_S)}
    &\le w,                                             \label{eq:inverse-series}\\
 \norm{\Tr\log(I+F)}_{\Bcal_S}
    &\le Dn[-\log(1-u)].                                 \label{eq:det-series}
\end{align}
For \eqref{eq:det-series}, each trace is bounded by $Dn$ times
the matrix Schur norm, and $\sum_{j\ge1}u^j/j=-\log(1-u)$.

The exact Gaussian density identity is
\begin{equation}
 p_M(x)=p_{M_0}(x)
 \exp\left\{-\frac12\Tr\log(I+F)
            -\frac12x^T(M^{-1}-M_0^{-1})x\right\}.
 \label{eq:fulldensity}
\end{equation}
It expands both the determinant and the inverse. It follows that
\begin{equation}
 \norm{p_M(x)}_{\Bcal_S}
 \le |p_{M_0}(x)|
       \exp\left\{\frac{Dn}{2}[-\log(1-u)]
                              +\frac w2\norm{x}^2\right\}.
 \label{eq:density-majorant}
\end{equation}
The quadratic estimate uses the row-and-column Schur bound in
\eqref{eq:inverse-series}, applied to the block Euclidean norms of $x$.

Equation \eqref{eq:deformedgap} gives
\[
 \Herm M_0^{-1}
       =M_0^{-*}(\Herm M_0)M_0^{-1}\ge V_*^{-1}I,
 \qquad |\det M_0|\ge(\lambda/2)^{Dn}.
\]
Using $u\le1/2$ and $w\le1/(2V_*)$, the right-hand side of
\eqref{eq:density-majorant} is bounded by
$(8V_*/\lambda)^{Dn/2}$ times the independent real Gaussian density
with covariance $2V_*I$.
For one $D$-dimensional block $G$ of that real Gaussian,
\begin{equation}
 \mathbb P(\norm{G}>R)
 \le e^{-R^2/(8V_*)}
            \mathbb E e^{\norm{G}^2/(8V_*)}
 =2^{D/2}e^{-R^2/(8V_*)}.
 \label{eq:tail}
\end{equation}

Set $q_b(x_b)=-\one_{\{\norm{x_b}>R\}}$.
Combining \eqref{eq:density-majorant} and \eqref{eq:tail} proves
the support-valued all-marked bound
\begin{equation}
 \left\|\int p_M(x)\prod_{b\in S}q_b(x_b)\,dx
       \right\|_{\Bcal_S}
 \le\widehat\tau_R^{\,n}.
 \label{eq:algebra-tail}
\end{equation}
All coefficient sums, integrals, and matrix series can be interchanged
using these absolute majorants. The indicators remain fixed measurable
multipliers. The same bound covers $n=1$, with no tree increments.

The inequality $v_1\le\lambda/4$ also ensures that evaluation of these
matrices retains accretivity on the required interpolation domains.
The determinant branches in \eqref{eq:fulldensity} agree with the
normalized Gaussian branch, by continuation from zero perturbation.

\section{Exact marked polymers and tree estimates}

For a nonempty $S$, define
\[
 f_S(s;A)=\int p_{C_S(s;A)}(x)\prod_{b\in S}q_b(x_b)\,dx,
 \qquad m(S)=f_S(1;A).
\]
These functions are holomorphic in a neighborhood of the partition
polytope. The BKAR formula \cite{AR,BKARnotes} gives a sum over forests.
At a forest mask the covariance matrix is block diagonal along forest
components. Derivatives in the forest edges remain internal to those
components. Therefore the forest sum factorizes exactly as
\begin{equation}
 m(S)=\sum_{\pi\text{ partition of }S}
                         \prod_{B\in\pi}\zeta(B),
 \label{eq:moment-partition}
\end{equation}
where $\zeta(\{b\})=m(\{b\})$ and
\begin{equation}
 \zeta(B)=\sum_{T\text{ tree on }B}\int_{[0,1]^T}
                          \partial_T f_B(s^T(t))\,dt
                  \qquad(|B|\ge2).
 \label{eq:zeta}
\end{equation}
This identity also holds coefficientwise in the formal source series.
Although covariance blocks can read overlapping background supports,
the masked Gaussian matrix itself is block diagonal.

Banach-valued Cauchy applied to \eqref{eq:algebra-tail}, in the
polydisc \eqref{eq:cauchyradii}, yields
\[
 \norm{\partial_Tf_S(s^T)}_{\Bcal_S}
 \le\widehat\tau_R^{\,n}\delta^{-(n-1)}
       \prod_{\{b,c\}\in T}\max(d_b,d_c)J_{bc}.
\]
Each parameter is differentiated once, so no factorial appears from
Cauchy. For integer $d\ge1$, $d^d\le e^d d!$, while
$\sum_bd_b=2n-2$ and $\max(d_b,d_c)\le d_bd_c$.
Thus
\begin{equation}
 \norm{\partial_Tf_S(s^T)}_{\Bcal_S}
 \le\widehat\tau_R^{\,n}(e^2/\delta)^{n-1}
                         \prod_{b\in S}d_b!\prod_{e\in T}J_e.
 \label{eq:tree-bound}
\end{equation}

For completeness, the exact degree-weighted tree count is
\begin{equation}
 \sum_{T\text{ tree on }[n]}\prod_{i=1}^n d_i!
       =(n-2)!\binom{3n-3}{n-2}.                         \label{eq:treecount}
\end{equation}
Indeed write $d_i=k_i+1$ with $\sum_i k_i=n-2$.
The Pruefer count for this degree sequence is
$(n-2)!/\prod_i k_i!$.
Multiplication by $\prod_i d_i!$ leaves
$(n-2)!\prod_i(k_i+1)$, whose sum is the coefficient of
$z^{n-2}$ in $(1-z)^{-2n}$.
In particular,
\begin{equation}
 \frac1{(n-1)!}\sum_T\prod_i d_i!
       =\frac1{n-1}\binom{3n-3}{n-2}\le8^{n-1}.           \label{eq:treecountbound}
\end{equation}

To sum marked sets of size $n$ containing a fixed root, label the other
$n-1$ vertices and divide by $(n-1)!$. Dropping distinctness only
increases the upper bound; set the edge majorant at coincident labels
to zero. For each fixed tree, successive leaf sums of
$e^{\kappa d(b,c)}J_{bc}$ are bounded by $A_\kappa^{n-1}$.
Equations \eqref{eq:tree-bound}--\eqref{eq:treecountbound} imply
\begin{equation}
 \sup_b\sum_{\substack{S\ni b\\|S|=n}}
             e^{\kappa\ell(S)}\norm{\zeta(S)}_{\Bcal_S}
           \le\widehat\tau_R^{\,n}B_0^{n-1}.
 \label{eq:anchored-n}
\end{equation}
The singleton bound is \eqref{eq:algebra-tail}.
Writing the canonical support representation as
$\zeta(S)=\sum_X\zeta(S;X)(A|_X)$, we obtain
\begin{equation}
 \sup_b\sum_{S\ni b,X}
 e^{\kappa\ell(S)+\kappa\ell_S(X)+\alpha|S|}
                       \norm{\zeta(S;X)}_\infty
 \le
 \frac{\widehat\tau_Re^\alpha}
      {1-B_0\widehat\tau_Re^\alpha}
 \label{eq:polymernorm}
\end{equation}
whenever the denominator is positive.

Expanding $\prod_b(1+q_b)$ and using \eqref{eq:moment-partition}
gives the exact hard-core gas
\begin{equation}
 Z_{\Lambda,R}
  =\sum_{\{S_1,\ldots,S_k\}\text{ pairwise disjoint}}
                            \prod_i\zeta(S_i).
 \label{eq:gas}
\end{equation}
The empty family has weight one. Compatibility concerns the
\emph{marked sets} $S_i$, not the enlarged background supports.
Replacing it by disjointness of the enlarged supports would change
the partition function.

\section{The logarithm and background-rooted bounds}

Use the scalar majorants
\[
 \widetilde z(S)=e^{\kappa\ell(S)+|S|}
                 \sum_Xe^{\kappa\ell_S(X)}
                               \norm{\zeta(S;X)}_\infty,
 \qquad a(S)=|S|.
\]
Equation \eqref{eq:polymernorm} with $\alpha=2$ and
\eqref{eq:cutoff-smallness} imply
\begin{equation}
 \sum_{S':S'\cap S\ne\varnothing}\widetilde z(S')e^{|S'|}
 \le |S|\,2e^2\widehat\tau_R\le |S|.
 \label{eq:kp}
\end{equation}
The hard-core polymer criterion and rooted cluster estimate
\cite{KP} give the absolutely convergent Mayer logarithm
\begin{equation}
 \log Z_{\Lambda,R}
  =\sum_{k\ge1}\frac1{k!}\sum_{S_1,\ldots,S_k}
           \varphi^T(S_1,\ldots,S_k)\prod_i\zeta(S_i).
 \label{eq:mayer}
\end{equation}
Here $\varphi^T$ is the sum over connected graphs on $[k]$ of
the product of $-\one_{\{S_i\cap S_j\ne\varnothing\}}$ over edges.
Repeated marked polymers are allowed and are incompatible with
themselves. The convergent expansion is the logarithm obtained by
turning on an auxiliary defect strength from zero; it also proves
zero-freeness on that defect polydisc.

Expand each $\zeta(S_i)$ in its support series and group by
\[
 M=\bigcup_iS_i,\qquad X=\bigcup_iX_i,
\]
obtaining $U(M;X)(A|_X)$. The set $M$ is always nonempty.
For a connected overlap graph of the marked polymers,
\begin{equation}
 \ell(M)\le\sum_i\ell(S_i),\qquad
 \ell_M(X)\le\sum_i\ell_{S_i}(X_i),\qquad
 |M|\le\sum_i|S_i|.
 \label{eq:unionweights}
\end{equation}
For the first inequality, combine the trees on the $S_i$;
overlapping components need no extra connecting edges.
For the second, a forest attached to $S_i$ is also attached to
the larger root set $M$, and \eqref{eq:forest-subadd} applies.
The scalar majorant and rooted cluster bound consequently imply
\begin{equation}
 \sup_b\sum_{M\ni b,X}
      e^{\kappa[\ell(M)+\ell_M(X)]+|M|}
                      \norm{U(M;X)}_\infty
      \le2e^2\widehat\tau_R.
 \label{eq:markedlogbound}
\end{equation}
To check the rooting, a cluster whose marked union contains $b$
has at least one polymer containing $b$. Choosing one such polymer
overcounts. Its rooted Mayer sum with coefficient $1/(k-1)!$
is at most $e^{a(S)}$. Summing its weighted activity uses exactly
\eqref{eq:polymernorm} at $\alpha=2$, so no uncontrolled factor
in the number of marks or polymers is introduced.

Choose deterministic minimizing metric trees on $M$ and minimizing
forests attaching $X$ to $M$, and realize their edges by deterministic
shortest lattice paths. Their union $H(M,X)$ is connected and contains
$M\cup X$. If
\[
 L(M,X)=\ell(M)+\ell_M(X),
\]
then
\begin{equation}
 |H(M,X)|\le|M|+L(M,X),\qquad
 d(a,b)\le L(M,X)\quad(a\in M,\ b\in H(M,X)).
 \label{eq:hull}
\end{equation}
If paths leave a nonconvex $\Lambda$, the coefficient still reads
only the original support $X\subset\Lambda$.

Group all $U(M;X)$ with the same $H(M,X)$ into $E_{\Lambda,H}$.
This is a strictly local holomorphic function of $A|_{H\cap\Lambda}$.
For $a\in M$, $b\in H$, and $a_*=\min\{1/2,\kappa/4\}$,
\begin{equation}
 e^{a_*|H|}
 \le e^{|M|+\kappa L(M,X)}e^{-(\kappa/2)d(a,b)}.
 \label{eq:rootcharge}
\end{equation}
Choose one marked $a$, or overcount by all $a\in M$.
Sum first over the clusters containing this $a$ using
\eqref{eq:markedlogbound}, and then over $a$.
Equation \eqref{eq:rootcharge} proves \eqref{eq:final-bound}.
This pays explicitly for a root lying in an added background support
or connecting path rather than in a marked block.

All series converge locally uniformly in the background and absolutely
in the displayed majorants. They therefore preserve holomorphy and
justify grouping and evaluation. This completes the proof
of Theorem \ref{thm:main}.

\section{Why no disconnected unmarked component remains}

The density expanded in \eqref{eq:fulldensity} is the normalized
\emph{marginal on a nonempty marked set}. Every Neumann path in
$\Delta C_S$ starts and ends in $S$.
Every background support in its coefficients is thus attached to $S$
before the inverse and determinant are expanded. Multiplication in
$\Bcal_S$ preserves this attachment property.

The factors $M_0^{-1}$ and the trace in the determinant series operate
only on the already designated roots $S$.
They cannot create a component consisting solely of remote
background supports. In the subsequent logarithm every gas polymer
contains at least one mark. There is no $M=\varnothing$ term.
This is the precise cancellation supplied by normalized marginals;
one need not begin with a bulk determinant and cancel disconnected
source-only loops afterward.

\section{Limits of the theorem and the model interface}

The theorem is a sufficient statement for a Gaussian cutoff
normalizer. Notebooks 06 and 07 now provide its positive, bounded,
exponentially local homogeneous flat reference on the specified
retained coordinates, including a free refinement limit at rate
$O(M^{-2})$ in weighted kernel norm. The physical identification has
the smooth-minimizer and coordinate requirements stated there.
The following points remain separate in a Yang--Mills application.
\begin{enumerate}
\item The actual independent-coordinate precision must admit
  \eqref{eq:precision} and the uniformly small norm
  \eqref{eq:local-norm} on the required complex domain.
  A spectral gap and covariance decay alone do not imply that
  local perturbation structure.
\item Gauge constraints and coordinate choices must lead to fixed
  block dimension and a fixed real cutoff domain during complex
  background continuation. A moving cutoff domain is not covered.
\item Baseline and perturbation constants must be uniform in scale,
  history, and volume, and fixed before the target remainder tolerance.
\item Nonquadratic interaction polymers require a joint expansion
  and stability argument. This notebook proves no such interacting
  transfer.
\item The positive exponent $a_*$ in \eqref{eq:final-bound} must be
  compared with the numerical polymer norm actually needed downstream.
  A stronger target exponent requires enough initial decay.
\item Strict background locality does not imply consistency between
  different finite volumes. The baseline $C_{0,\Lambda}$ can change,
  affecting $E_{\Lambda,H}$ even for fixed $H$.
  Likewise, if a baseline is allowed to depend on an external real
  reference configuration, locality in the perturbation variables
  does not establish locality in that reference configuration.
  Applying the theorem around each reference separately cannot remove
  this additional dependence.
\item A weak gauge-potential ball or constrained current domain is
  not silently identified with the product domain used here.
  An appropriate chart or a separate support-algebra argument is needed.
\item The theorem does not assert activity-level gauge or lattice
  rotation invariance. Deterministic connecting paths can preserve
  translations while breaking rotations. Any equivariant construction
  or averaging step must preserve the actual domain, support labels,
  and weighted norm. Gauge invariance of the total Gaussian functional
  alone does not establish gauge invariance of each activity.
\end{enumerate}

Covariance entry decay alone cannot replace the locality hypothesis.
For example,
\[
 C_{bc}(A)=(1+\varepsilon A_{x_*})I_D\,\one_{\{b=c\}},
 \qquad |A_{x_*}|<1,\quad 0<\varepsilon<1/2,
\]
is uniformly accretive, bounded, holomorphic, and has zero
off-diagonal covariance. Yet every singleton cutoff logarithm
depends on the arbitrarily remote background coordinate $A_{x_*}$.
The local precision hypothesis excludes precisely this type of
uncontrolled remote dependence.

Finally, if $R(g)=c_0(\log g^{-2})^p$ with $p>1/2$ and $c_0>0$,
then the right-hand side of \eqref{eq:final-bound} is a fixed
constant times
$\exp[-c_0^2(\log g^{-2})^{2p}/(8V_*)]$.
It is smaller than any prescribed fixed multiple of
$g(\log g^{-2})^q$ for sufficiently small $g$, for every fixed
real $q$. The threshold is independent of a remainder tolerance
only when all preceding constants were chosen independently of
that tolerance.

\begin{thebibliography}{9}
\bibitem{AR}
A. Abdesselam and V. Rivasseau,
\emph{Trees, forests and jungles: a botanical garden for cluster expansions},
\href{https://arxiv.org/abs/hep-th/9409094}{arXiv:hep-th/9409094}.

\bibitem{BKARnotes}
A. Abdesselam,
\emph{Notes on the Brydges--Kennedy--Abdesselam--Rivasseau forest
interpolation formula}, Theorem 1 and the partition-polytope construction,
\href{https://mabdesselam.github.io/BKAR.pdf}{author's notes}.

\bibitem{KP}
R. Koteck\'y and D. Preiss,
\emph{Cluster expansion for abstract polymer models},
Communications in Mathematical Physics \textbf{103} (1986), 491--498,
\href{https://doi.org/10.1007/BF01211762}{doi:10.1007/BF01211762}.
\end{thebibliography}

\end{document}
